% Imporditud: tools/import_eksperimendid.py
% Allikas: Eesti füüsikaolümpiaadi eksperimendi ülesannete
%   kogu (Rael Kalda, 2023), ülesanne Ü144, lahendus L144.
\setAuthor{Jaan Kalda}
\setRound{piirkonnavoor}
\setYear{2014}
\setNumber{G E2}
\setDifficulty{4}
\setTopic{staatika}
\setVahendid{Kümme spagetikõrt, teadaoleva massiga joonlaud}
\setPunktid{12}
\setAllikas{Eksperimendi kogu Ü144, lahendus lk 107}
\setLahendusseis{toores}

\prob{Spagetid}
Leidke spageti murdumiseks vajalik jõumoment murdumispunkti suhtes.

\hint

\solu
Lahendamise esialgne idee: fikseerime spageti ühe otsa, nt hoiame sealt näpuga

kinni, ja rakendame teise otsa järjest kasvavat jõudu seni kuni spageti murdub.

Mõõtes jõu õla (kauguse jõu rakenduspnktist kuni murdumiskohani --- mis on hari-

likult otse näpu juures, maksimaalsel kaugusel jõu rakenduspunktist) ning teades

rakendatud jõu väärtust saame leida spageti äramurdunud otsale mõjunud jõumo-

mendi murdumispunkti suhtes. Enne murdumist oli see välise jõu moment taskaa-

lustatud spageti ristlõikes mõjunud pingejõudude momendiga: see ongi suurus,

mida otsime. Tasakaalutingimuse tõttu pidi see otsitav suurus olema võrdne meie

poolt rakendatud jõu momendiga.

Katse käigus selgub, et tüüpiline spageti on nii painduv, et täispikkuses kõrre pu-

hul on murdumishetkel kõrre otste vaheline nurk liiga suur: jõu õlga on raske

mõõta. Seepärast on mugavam mõõta lühema kõrrejupi murudmist, nt hoides kin-

ni kõrre keskpunktist. Selgub, et see idee klapib hästi joonlaua kaaluga: kui pool

joonlaua kaalust mg (kus m on joonlaua mass) rakendada kõrre otsale (selleks lase-

me peaaegu horisontaalsel joonlaual toetuda ühe otsaga vastu lauda ning toetame

teist otsa spagetiga), siis kriitiliseks spageti õlapikkuseks (mille juures see mur-

dub) ongi umbes pool spageti pikkusest. Kasutame alguses lühemat õlga nii, et

spageti veel ei murdu ning suurendame õlga kuni murdumiseni; mõõdame mur-

dunud spagetijupi pikkuse l (kui joonlaud ei toetunud rangelt spageti otspunkti-

le, siis tuleb lahutada tulemist toetuspunkti kaugus otspunktist (eeldatavasti paar

millimeetrit). Kui kõrre kõverus ei olnud murdumishetkel väga suur, siis ongi l jõu-

õlaks ning otsitav jõumoment M = mgl/2. Kui kõrre kõverus oli märkimisväärne,

siis tuleks täpsema tulemuse huvides mõõta joonlaua abil kõrre kõõlu pikkus L

murdumishetkel (sedasama murdunud kõrrejuppi uuesti painutades).

Teeme korduskatseid ning arvutame mõõtmistulemuste keskmise. Hea oleks kasu-

tada erinevaid jõudusid: eelpoolkirjeldatud viisil rakendasime pool joonlaua kaa-

lust, kuid kui tõsta joonlauda spagetiga joonlaua keskpunkti lähedalt, siis on või-

malik saada ka joonlaua täiskaal mg (või ka nt 2 mg, kui spageti toetuspunkt on joonlaua otsast veerandi joonlaua pikkuse kaugusel).

Hindamine: Idee arvutada jõumoment kui spageti otsale rakendatud jõu ning selle jõu õla (murdumispunkti suhtes) korrutis (3 p). Idee määrata jõu õlg murudunud spagetiotsa abil (1 p) ning reaalse mõõtmise teostamine (1 p). Valemi M = mgl/2 (või vastavalt

olukorrale õige avaldise, nt M = mgl või M = 2 mgl) rakendamine (2 p) ja õiges suurusjärgus vastuse saamine (1 p).

Kordusmõõtmiste tegemine: iga lisamõõtmine kuni viienda mõõtmistulemseni: a

0,5 punkti. Seega lisamõõtmiste eest täiendavalt 2 p.

Vastus, mis ei erine õigest rohkem kui 50\%: 0,5 punkti ning kui see ei erine rohkem

kui 30\%, siis kokku --- 1 p.

Mainitakse, et spagetikõrre kõverus on nii väike, et spagetitüki pikkus on peaae-

gu võrdne jõu õlaga või kui on mõõdetud spagetitüki pikkuse asemel selle kõõlu

pikkus painutatud olekus --- 1 p.
\probend
